% ID: SEL-D0025-OA-E01
% ESTADO: validada. S026, cotejada con examen y criterios.
Sean $x$ los litros de leche entera e $y$ los de desnatada.
El modelo es $x,y\ge0$, $x\le2y$, $x+y\le3000$ y $x\le1200$.
El beneficio diario es $G=0.4x+0.1y$.

Para $0\le x\le1200$, los límites son $x/2\le y\le3000-x$.
El recinto es el cuadrilátero con vértices $(0,0)$, $(1200,600)$,
$(1200,1800)$ y $(0,3000)$.
\begin{center}
\setlength{\unitlength}{2mm}
\begin{picture}(22,33)(-4,-1)
\qbezier(0,0)(9,0)(18,0)\put(18.2,0){$x$}
\qbezier(0,0)(0,15)(0,30)\put(0,30.6){$y$}
\qbezier(0,0)(6,3)(12,6)
\qbezier(12,6)(12,12)(12,18)
\qbezier(12,18)(6,24)(0,30)
\put(0.5,29){$(0,3000)$}
\put(12.3,18){$(1200,1800)$}
\put(12.3,6){$(1200,600)$}
\end{picture}
\end{center}
En esos vértices el beneficio vale, respectivamente,
$0$, $540$, $660$ y $300$ euros. Como la función es lineal,
el máximo se alcanza en el vértice de mayor valor:
\textbf{1200 litros de entera y 1800 de desnatada}, con
$\boxed{660}$ euros de beneficio diario.
